% Generated by docs/tools/params.py from the X30 sources: do not edit.
\subsubsection*{\code{x30_driver}}
\begin{small}
\begin{longtable}{@{}>{\raggedright}p{0.27\linewidth}>{\raggedright}p{0.15\linewidth}>{\raggedright}p{0.08\linewidth}>{\raggedright}p{0.1\linewidth}>{\raggedright\arraybackslash}p{0.32\linewidth}@{}}
\toprule
\textbf{Parameter} & \textbf{Default} & \textbf{Unit} & \textbf{Range} & \textbf{Description} \\
\midrule
\endhead
\bottomrule
\endfoot
\code{motion_host_ip} & \code{192.168.1.103} & & & network.toml (X30 Application Manual 4.3.2; README "Network"). \\
\code{motion_host_port} & \code{43893} & & &  \\
\code{local_ip} & \code{0.0.0.0} & & & address the report socket binds to \\
\code{local_port} & \code{43897} & & & the port listed in network.toml \\
\code{heartbeat_hz} & \code{5.0} & & & 1.2.1: at least 2 Hz (lower values: 2), or the link drops \\
\code{connect_retry_sec} & \code{5.0} & & & resend the connect command this often while no reports come \\
\code{command_prefix} & \code{0x31} & & & automatic mode switch and the takeover check follow the prefix. \\
\code{save_data_code} & \code{0x010C01} & & & emergency stop or fall: the robot exits its motion program. \\
\code{perception_host_ip} & \code{192.168.1.105} & & & Perception host (1.4 autonomous charging, 1.5 elevation map). Empty: no perception services. \\
\code{elevation_map_port} & \code{43899} & & & 1.5: commands and the 10 Hz mode feedback \\
\code{charge_port} & \code{3333} & & & 1.4: request / response \\
\code{lio_port} & \code{60000} & & & 1.5.2.4/5: LiDAR-inertial odometry on / off (multi-frame mode) \\
\code{obstacle_height_code} & \code{0x3100EE04} & & & perception host does not change, try 0x3101EE04 (the other commands are 0x3101EE0x). \\
\code{elevation_query_sec} & \code{5.0} & & & restart the elevation map feedback this often when it stops \\
\code{charge_query_sec} & \code{2.0} & & & ask for the charging state this often (0: never) \\
\code{feedback_timeout_sec} & \code{2.0} & & & how long a service waits for the robot to report the change \\
\code{velocity_mode} & \code{axis} & & & motion program runs /cmd\_vel\_corrected in non-manual mode, 2.3). \\
\code{velocity_topic} & \code{/cmd_vel_corrected} & & &  \\
\code{cmd_rate_hz} & \code{50.0} & & & 1.2.3: 50 Hz; the robot drops an axis after 1 s \\
\code{cmd_timeout_sec} & \code{0.5} & & & cmd\_vel older than this: zeros \\
\code{stop_burst_sec} & \code{0.5} & & & zeros streamed after motion (at least one), then silence \\
\code{max_vx} & \code{1.0} & & & m/s, clamp of cmd\_vel; the gait's own maximum applies too \\
\code{max_vy} & \code{0.5} & & & m/s \\
\code{max_wz} & \code{1.0} & & & rad/s \\
\code{axis_dead_zone} & \code{655} & & & 1.2.3.2: |value| < 655 is 0; smaller commands go out as 656 \\
\code{gait_limits} & \code{{}} & & & number; listed Appendix B gaits are replaced. Example: \{silent: [0.5, 0.5, 0.2, 0.5]\} \\
\code{invert_vy} & \code{false} & & & true if +linear.y moves the robot RIGHT (README checklist) \\
\code{invert_wz} & \code{false} & & & true if +angular.z turns the robot RIGHT (README checklist) \\
\code{auto_control_mode} & \code{true} & & & episode start: that mode (fresh report), 3 tries in 2 s \\
\code{auto_stepping} & \code{true} & & & episode start: stepping from torque standing, 3 tries in \\
\code{idle_stop_stepping_sec} & \code{0.0} & & & stop stepping after this long without motion (0: never) \\
\code{auto_torque_control} & \code{true} & & & stand\_up continues from initial standing to torque control \\
\code{stand_timeout_sec} & \code{10.0} & & & stand\_up / sit\_down wait this long for the robot \\
\code{transition_timeout_sec} & \code{5.0} & & & torque control, stepping and e-stop wait this long; also how \\
\code{odom_frame} & \code{odom} & & &  \\
\code{base_frame} & \code{base_link} & & &  \\
\code{imu_frame} & \code{imu_link} & & &  \\
\code{publish_tf} & \code{false} & & & odom -> base\_link on /tf (off: the odometry stack owns it) \\
\code{odom_hz} & \code{50.0} & & & one odom per new motion state report, at most this rate \\
\code{imu_hz} & \code{100.0} & & & one imu per new sensor report, at most this rate \\
\code{joint_state_hz} & \code{50.0} & & &  \\
\code{state_hz} & \code{10.0} & & & state (diagnostic\_msgs/DiagnosticStatus) \\
\code{battery_hz} & \code{1.0} & & & the robot reports its battery at 0.5 Hz \\
\code{state_timeout_sec} & \code{1.0} & & & no report for this long: disconnected (0x1009: no axes) \\
\code{low_battery} & \code{0.2} & & & 0-1, the state warns below \\
\code{battery_voltage_scale} & \code{0.0} & & & V per unit; 0: V, or mV / 10 mV recognised by size \\
\code{battery_current_scale} & \code{0.01} & & & A per unit: 0.01 (10 mA, Deeprcs >= 1.10.19) or 1.0 \\
\end{longtable}
\end{small}
